\section{QUICKPRUNE: A PRUNING ALGORITHM WITH GUARANTEES} \label{sec:alg} 
In this section, we describe the pruning algorithms introduced in this paper.
In Section \ref{sec:prune-single},
we describe the pruning algorithm for a single knapsack constraint.
The main pruning algorithm (\qs, Alg. \ref{alg:multi}) runs multiple
copies of the single constraint algorithm in parallel, as detailed
in Section \ref{sec:prune-multi}.

To prove theoretical guarantees for \qs, we make
the following assumption on the costs. Intuitively, the assumption
says that no element of an optimal solution consumes a very large
fraction of the total budget. 
\begin{assumption}[No Huge Items (\nhi)]
  Given an instance $(f,c,\kappa)$ of \sm,
  the instance satisfies the assumption with $\eta > 0$
  if there exists an optimal solution $O$ to the
  instance, such that, for all $o \in O$, $c(o) \le \kappa (1-\eta)$.
\end{assumption}
For example, in the case of size constraint, this assumption is
satisfied if $\kappa \ge 2$ and $\eta \le 1/2$. 

Our theoretical guarantees for the main pruning algorithm are
summarized in Theorem \ref{thm:multi}, which is proved in Section \ref{sec:prune-multi}.
Observe that we show a constant factor $C$ of the optimal value is retained
for all budgets in the range $[\kappa_{\min}, \kappa_{\max}]$; the desired
range of budgets to support by the user. The constant
depends mainly on $\gamma$, the submodularity parameter of $f$,
and the input parameter $\delta$; $C$ is optimized
with $\delta = 1$ at the value $1/24 - \epsi$.
Also, we show a size bound on the size of $\uni'$, the pruned ground set,
which depends logarithmically on the size $n$ of $\uni$, as well as the values
of the maximum and minimum budgets and the minimum cost of an element. 
\begin{theorem} \label{thm:multi}
  Let $\kappa_{\min} < \kappa_{\max}$, let $0 < \eta \le 1/2$, and let $f, c$ be
  $\gamma$-submodular and modular functions, respectively.
  Suppose that for all $\kappa' \in [\kappa_{\min}, \kappa_{\max}]$,
  the instance $(f,c,\kappa')$ satisfies the \nhi assumption with $\eta$.
  Let $\uni'$ be the output of \qs (Alg. \ref{alg:multi})
  with parameters
  $(f,c,\kappa_{\min},\kappa_{\max}, \delta, \epsi, \eta)$.
  Then, for all $\kappa' \in [\kappa_{\min}, \kappa_{\max}]$,
  it holds that
  $f_{\kappa'}( \uni' ) \ge
  \frac{\delta \gamma^5(1 - \epsi \gamma^{-1})}{6(\delta \gamma^2 + 1)( 1 + \gamma^{-1} \delta ) }  f_{\kappa'}( \uni ).$
  And, $| \uni' | = \oh{ \log( \frac{ \kappa_{\max} }{ \kappa_{\min} } ) \left( 1 + \frac{\kappa_{\max} }{\delta c_{min} } \right) \log ( n / \epsi ) }.$
\end{theorem}
\subsection{Pruning for a Single Knapsack Constraint} \label{sec:prune-single}
First, we target the case of a single constraint.
Given a single budget value $\kappa$, we formulate an algorithm to prune for the
instance $(f, \uni, \kappa,c)$. That is, we produce an algorithm that, for some
$\alpha > 0$, produces $\uni'$ such that $f_{\kappa}( \uni' ) \ge \alpha  f_{\kappa }( \uni )$
and $|\uni'| \le \oh{F( \kappa, c) \polylog(n)}. $

The algorithm \qss is given in Alg. \ref{alg:single}. In addition to the problem
instance, it takes parameter $\delta > 0$, which impacts the size of the resulting
pruned set $\uni'$ by adjusting the condition to add elements to $\uni'$;
and parameter $\epsi > 0$, which controls how aggressively the algorithm deletes elements.
In overview, the algorithm works by taking one pass through the ground set $\uni$ and
processing elements one-by-one. An element is added to $\uni'$ if it meets the condition
on Line \ref{line:add}: $\marge{e}{A} \ge \frac{ \delta c(e) f(A) }{ \kappa }$; intuitively,
this means the element should increase the value of $A$ (the pruned set)
an amount proportional to $c(e) / \kappa$ to be worth retaining.

To ensure the size of $\uni'$ doesn't grow too large, a deletion condition is checked
on Line \ref{line:delete}, which intuitively asks if the value of $A$ has increased by
a large factor from the previous checkpoint; if so, the original elements of $A$ are
discarded. Submodularity is used to bound the amount of value lost from the deletion,
and the condition on element addition bounds the maximal number of elements required for
the increased value.

The theoretical properties are summarized in the following theorem, proven in Section \ref{sec:single-proofs}.
\begin{theorem} \label{thm:single}
  Let Alg. \ref{alg:single} be run on instance $(\uni, c, f, \kappa)$, with parameters $\epsi, \delta > 0,$
  such that $f$ is $\gamma$-submodular. 
  Then, Alg. \ref{alg:single} produces pruned universe $\uni'$, such that
  $| \uni ' | < 2 \left( 1 + \frac{\kappa }{\delta c_{min} } \right) \log ( n / \epsi ) + 3,$
  and there exists $A' \subseteq \uni'$, such that $c(A') \le \kappa$ and
  $f(A') \ge \frac{\delta \gamma^4(1 - \epsi \gamma^{-1})}{2(\delta \gamma^2 + 1)( 1 + \gamma^{-1} \delta ) } \opt$.
  %$f(A') \ge \frac{\delta (1 - \epsi)}{2(\delta + 1)^2} \opt$.
  %Maximizing the ratio over $\delta$ yields $f(A') \ge \frac{(1 - \epsi)}{8} \opt$ at $\delta = 1$. 
\end{theorem}
\begin{algorithm}[t]
	\caption{\qss:  Pruning for a single constraint.}
	\begin{algorithmic}[1] \label{alg:single}
		\STATE \textbf{Input:} Instance $(f, \uni, c, \kappa)$, size-control parameter $\delta > 0$, deletion parameter $\epsi > 0$
		\STATE \textbf{Output:} Pruned ground set, $\uni'$
		\STATE \textbf{Initialize:} $A \gets \emptyset$, $a^* \gets \emptyset$, $A_s \gets \emptyset$
                \FOR {$e \in \uni$}
                \IF{ $c(e) > \kappa$  }
                \STATE \textbf{continue}
                \ENDIF{}
                \IF{$ \marge{e}{A} \ge \frac{\delta c(e) f(A) }{ \kappa }$ \label{line:add}}
                \STATE $A \gets A + e$
                \ENDIF{}
                \IF{ $f(e) > f(a^*)$  }
                \STATE $a^* \gets e$
                \ENDIF
                \IF{ $f(A) > \frac{ n }{ \epsi } f(A_s) $\label{line:delete} }
                \STATE $A \gets A \setminus A_s$ \label{line:delete-set}
                \STATE $A_s \gets A$
                \ENDIF
                \ENDFOR
                \STATE \textbf{return} $\uni' \gets A + a^*$ 
	\end{algorithmic}
\end{algorithm}
\subsubsection{Proof of Theorem \ref{thm:single}} \label{sec:single-proofs}
Due to space constraints, omitted proofs are provided in
Appendix \ref{apx:proof}. We require the following fact about geometrically increasing sequences of real
numbers. For our purposes, the sequence $(y_i)$ will assume the values of $f(A)$ as elements are
added. 
\begin{fact} \label{fact:1}
  Let $(y_i)_{i=1}^m$ be a sequence of positive real numbers, such that, for some $\beta > 0$, it holds that $y_i \ge (1 + \beta)y_{i-1}$, for all $i \in [m]$. Let $\gamma > 0$. Then, if $m \ge \frac{\beta + 1}{\beta}\log \gamma^{-1} $, it holds that $y_m \ge y_1 / \gamma$.
\end{fact}
%\begin{proof} ...
%\end{proof}
Using Fact \ref{fact:1}, we bound the number of elements in $\uni'$. Intuitively, the geometric increase
in the value of $f(A)$ yields a bound on the maximum number of elements until the deletion condition
on Line \ref{line:delete} is triggered. 
\begin{proposition} \label{prop:uprime-size}
  The size of $\uni'$, as returned by Alg. \ref{alg:single}, satisfies
  $|\uni'| < 2 \left( 1 + \frac{\kappa }{\delta c_{min} } \right) \log ( n / \epsi ) + 3,$
  where $c_{min} = \min_{a \in \uni } c(a).$
\end{proposition}
Next, we bound how much value may have been lost from deletion throughout the execution
of the algorithm. In the following proposition, $\hat A$ is all elements ever added to
$A$, and $\dot A$ is the actual set obtained by the algorithm after all deletions. We
show at most an $\epsi\gamma^{-1}$-fraction of value is lost. 
\begin{proposition} \label{prop:ahat-adot}
  Suppose $f$ is $\gamma$-submodular.
  Let $A_i$ be the value of $A$ after the execution of iteration $i$ of the \textbf{for} loop,
  for $i = 1$ to $n$, and $A_0 = \emptyset$. 
  Let $\hat A = \bigcup A_j, \dot A = A_n$. Then $f( \dot A ) \ge (1 - \gamma^{-1} \epsi )f( \hat A)$. 
\end{proposition}
So far, we have established a size bound on $\uni'$ and bounded the value of
$\uni'$ lost from deletion. Next, we turn to showing that there exists a feasible set
inside $\uni'$ that has a constant fraction of the optimal value.
We start by showing in Lemma \ref{lemma:value-inside-A}
and Proposition \ref{prop:aprime} that
a set $A^*$ exists within $\uni'$ that has a constant fraction
of $f(\uni')$. 
\begin{lemma} \label{lemma:value-inside-A}
  Let $f$ be $\gamma$-submodular.
  Let $\delta, \kappa > 0$. Suppose $A_i = \{ a_1, \ldots a_i \}$ is a sequence of sets satisfying
  (1) $c(a_i) \le \kappa$, and (2) $\marge{ a_i }{ A_{i - 1} } \ge \gamma \frac{\delta c( a_i ) f( A_{i -1 } ) }{ \kappa }$,
  for each $i$ from $1$ to $m$. Let $f( a^* ) \ge \max_{i \in [m]} f(a_i)$, and $c(a^*) \le \kappa$.
  Then there exists $A^* \subseteq A_m + a^*$, such that $c( A^* ) \le \kappa$, and
  $f( A^* ) \ge \frac{ \delta \gamma^4}{2(\delta\gamma^2 + 1)} f(A_m)$.
\end{lemma}
\begin{proposition} \label{prop:aprime}
  Let $\uni' = \dot A + a^*$ have its value at termination of \qss.
  There exists $A' \subseteq \uni' = \dot A + a^*$, such that  $f(A') \ge \frac{\delta \gamma^4}{2(\delta \gamma^2 + 1)} f(\dot A )$
  and $c(A') \le \kappa$. 
\end{proposition}
Next, we relate the value of $f(\uni')$ to
$\opt = f_\kappa ( \uni )$ in Proposition \ref{prop:opt-ahat}.
\begin{proposition} \label{prop:opt-ahat}
  Suppose Alg. \ref{alg:single} is run on instance $(\uni,c,f,\kappa )$ with
  parameters $\delta, \epsi > 0$. 
  Let $\hat A = \bigcup_{j \in [n]} A_j$ be all elements added to $A$.
  Then $f( \hat A ) \ge \opt / (1 + \gamma^{-1}\delta )$, where $\opt = \max_{S \subseteq \uni: c(S) \le \kappa } f(S).$
\end{proposition}
\begin{proof}[Proof of Theorem \ref{thm:single}]
  Finally, we are ready to put all the pieces together.
  Assume the hypotheses of the theorem statement. 
  The size bound on $\uni'$ follows by Proposition \ref{prop:uprime-size}.
  Let $A_j$ be the value of the set
  $A$ at the beginning of the $j$th iteration of \qss.
  Let $\dot A = A_{n+1}$ be the final value of $A$,
  $\hat A = \bigcup_{i \in [n + 1]} A_i$.
  Let $A'$ be as guaranteed by Proposition \ref{prop:aprime}.
  Then$ f(A') \overset{(a)}{\ge} \frac{\delta \gamma^4 }{2(\delta \gamma^2 + 1)} f( \dot A ) 
  \overset{(b)}{\ge} \frac{\delta \gamma^4(1 - \epsi \gamma^{-1})}{2(\delta \gamma^2 + 1)} f( \hat A )$,
  %        &\overset{(c)}{\ge} \frac{\delta \gamma^4(1 - \epsi \gamma^{-1})}{2(\delta \gamma^2 + 1)( 1 + \gamma^{-1} \delta ) } \opt,
  where (a) is by Prop. \ref{prop:aprime},
  (b) is by Prop. \ref{prop:ahat-adot},
  and the result follows from Prop. \ref{prop:opt-ahat}.
\end{proof}

\subsection{The Pruning Algorithm: \qs} \label{sec:prune-multi} 
\begin{algorithm}[t]
	\caption{\qs: The Pruning Algorithm.}
	\begin{algorithmic}[1] \label{alg:multi}
          \STATE \textbf{Input:} Instances $(f, \uni, c, [\kappa_{\min}, \kappa_{\max}])$,  parameters $\delta > 0$, $\epsi > 0$,
          $0 < \eta \le 1/2$.
		\STATE \textbf{Output:} Pruned ground set, $\uni'$
		\STATE $\mathcal B = \{ \tau_i = \kappa_{\max}(1 - \eta )^i : i \in \mathbb Z, (1 - \eta) \kappa_{\min} \le \tau_i \le \kappa_{\max}\}$\label{line:multi-budgets}
                \STATE Initialize a copy of $\mathcal QP_{\tau}$ of \qss for each $\tau \in \mathcal B$, with parameters $\epsi, \delta$.
                \FOR {$e \in \uni$}
                \STATE Pass $e$ to $\mathcal QP_{\tau}$ for all $\tau \in \mathcal B$.
                \ENDFOR
                \STATE \textbf{return} the union of all sets returned by all instances $\mathcal QP_{\tau}$
	\end{algorithmic}
\end{algorithm}
In this section, we describe the main pruning algorithm \qs (Alg. \ref{alg:multi})
and prove Theorem \ref{thm:multi}. 

The algorithm \qs works as follows. As input, it takes instances with a range of
budgets $[\kappa_{\min}, \kappa_{\max}]$, the parameters $\epsi, \delta$ for
\qss, and a parameter $\eta > 0,$ such that all of the instances satisfy
Assumption \nhi with $\eta$. The algorithm then runs $\log ( \kappa_{\max} / \kappa_{\min} )$
copies of \qss in parallel -- one each for budget $\tau_i = \kappa_{max}(1 - \eta )^i$ in the
range $[( 1 - \eta ) \kappa_{\min}, \kappa_{\max} ]$. It returns the union of all pruned sets
returned from each copy of \qss.

To establish Theorem \ref{thm:multi}, we first show the following proposition,
with implications to how Assumption \nhi can relate an optimal solution of a given
budget to one of the instances handled by one of the copies of \qss. 
\begin{proposition} \label{prop:part}
  Suppose instance $(f,c,\kappa)$ satisfies the \nhi assumption
  with $0 < \eta \le 1/2$. 
  Then, there exists an optimal solution $O$ to the instance, such
  that $O$ can be partitioned into at most three sets $\{ O_i : i \in [3] \}$, such
  that $c(O_i) \le \kappa (1 - \eta )$, for each $i \in [3]$. 
\end{proposition}
\begin{proof}[Proof of Theorem \ref{thm:multi}]
  Assume the hypotheses of the theorem, and let
  $\kappa' \in [\kappa_{\min}, \kappa_{\max}]$. By the choice of
  the set $\mathcal B$ on Line \ref{line:multi-budgets}, there
  exists a $\tau \in \mathcal B$, such that $\kappa'(1 - \eta) \le \tau \le \kappa'$.
  Let $O \subseteq \uni$ be an optimal solution to $f_{\kappa'}( \uni )$
  satisfying Assumption \nhi. By Proposition \ref{prop:part}, $O$ can be partitioned
  into at most three sets $\{ O_i \}$, such that $c( O_i ) \le \kappa ( 1 - \eta ) \le \tau$.

  Moreover, by Theorem \ref{thm:single}, there exists $X \subseteq \uni'$, such that
  1) $c(X) \le \tau$, and 
  2) $f(X) \ge
  \alpha f( O_i )$,
  with $\alpha =  \frac{\delta \gamma^4(1 - \epsi \gamma^{-1})}{2(\delta \gamma^2 + 1)( 1 + \gamma^{-1} \delta ) }$.
  Therefore, $3f(X) \ge \alpha \sum_{i = 1}^3 f( O_i ) \ge \alpha \gamma f(O)$, where the last inequality follows from submodularity.
  Thus, $f(X) \ge \alpha \gamma f(O) / 3$.
  Finally, the size bound on $\uni'$
  follows from Prop. \ref{prop:uprime-size} and the size of $\mathcal B$. 
\end{proof}
% \begin{algorithm}[htb]
% 	\caption{Sparsification for Size Constraint}
% 	\begin{algorithmic}[1]\label{alg:size_constraint}
% 		\STATE \textbf{Input:} ground set $V$, oracle $f$,  $\delta$, size constraint $k$
% 		\STATE \textbf{Output:} pruned ground set $V'$
% 		\STATE \textbf{Initialize:} $V'\leftarrow\emptyset$
%             \FOR {$v\in V$}
%                 % \STATE Remove from $V$ the top $(1-1/\sqrt{c})|V|$ of elements with the smallest $w_{Uv}$;
%                 \IF{$f(V'+v)-f(V') \geq \frac{\delta f(V')}{k}$}
%                 \STATE $V'\leftarrow V'\cup v$
%                 \ENDIF{}
%             \ENDFOR
  
%   % $n\leftarrow|V|$
% % 		\WHILE{$|V|>r\log n$}
% % 		\STATE Sample $r\log n$ items uniformly at random from $V$ and place them in $U$;
% % 		\STATE $V\leftarrow V\backslash U$;
% % 		\STATE $V'\leftarrow V' \cup U$;
% % %		\FOR {$v\in V$, $u\in R$}
% % %		\IF{$f(v|u+S)\leq\epsilon$}\STATE{$V:=V\backslash v$}\ENDIF
% % %		\ENDFOR
% % 		\FOR {$v\in V$}
% %                 % \STATE Remove from $V$ the top $(1-1/\sqrt{c})|V|$ of elements with the smallest $w_{Uv}$;
% %                 \IF{$f(V'+v)-f(V') \geq \frac{\delta f(V')}{k}$}
% %                 \STATE $V'\leftarrow V'\cup v'$
% %                 \ENDIF{}
% %             \ENDFOR
% % 		\STATE Remove from $V$ the top $(1-1/\sqrt{c})|V|$ of elements with the smallest $w_{Uv}$;
% % 		\ENDWHILE
% % 		\STATE $V'\leftarrow V\cup V'$
% 	\end{algorithmic}
% \end{algorithm}

% \begin{algorithm}[htb]
% 	\caption{Sparsification for knapsack Constraint}
% 	\begin{algorithmic}[1]\label{alg:knapsack_constraint}
% 		\STATE \textbf{Input:} ground set $V$, oracle $f$, $\delta$,  $k$ , $w$
% 		\STATE \textbf{Output:} $V'$
% 		\STATE \textbf{Initialize:} $V'\leftarrow\emptyset$
%             \FOR {$v\in V$}
%                 % \STATE Remove from $V$ the top $(1-1/\sqrt{c})|V|$ of elements with the smallest $w_{Uv}$;
%                 \IF{$ \frac{f(V'+v)-f(V') }{w(v)} \geq \frac{\delta f(V')}{k}$}
%                 \STATE $V'\leftarrow V'\cup v$
%                 \ENDIF{}
%             \ENDFOR
  

% 	\end{algorithmic}
% \end{algorithm}

% \begin{algorithm}[htb]
% 	\caption{Sparsification for Multiple Budget for size constraint}
% 	\begin{algorithmic}[1]\label{alg:knapsack_constraint}
% 		\STATE \textbf{Input:} ground set $U$, oracle $f$, $\epsilon$, size constraint $k$ , $\delta$
% 		\STATE \textbf{Output:} $U'$
%             \STATE \textbf{Initialize:} $U' \leftarrow \emptyset $ 
% 		\STATE $ m \leftarrow \left\lceil \frac{log(k)}{log(1+\epsilon)} \right\rceil$ 

%             \FOR {$i \leftarrow 0$ to $m$} 
%                 \STATE $\tau_{i} \leftarrow (1+ \epsilon)^i$ , $U_{\tau_{i}} \leftarrow \emptyset$
%             \ENDFOR

%             \FOR {$v \in U$} 
%                 % \STATE $\tau \leftarrow (1+ \epsilon)^i$ , $U_\tau \leftarrow \emptyset$
%                 \FOR{$i \leftarrow 0$ to $m$}
%                     \IF{$f(U_{\tau_{i}}+v)-f(U_{\tau_{i}}) \geq \frac{\delta f(U_{\tau_{i}})}{\tau_{i}}$}
%                     \STATE $U_{\tau_{i}}  \leftarrow U_{\tau_{i}} \cup v$
%                 \ENDIF{}
%                 \ENDFOR
                
%             \ENDFOR
%             \STATE $U' \leftarrow \cup_{i \in [m]} U_{\tau_{i}} $
%             % \FOR {$v \in U$} 
                
%             %     \FOR{$i \leftarrow 0$ to $m$ }
%             %         % \STATE $\tau \leftarrow (1+ \epsilon)^i$ , $U_\tau \leftarrow \emptyset$
%             %         \IF{$f(U_\tau+v)-f(U_\tau) \geq \frac{\delta f(U_\tau)}{\tau}$}
%             %         \STATE $U_\tau\leftarrow U_\tau\cup v$
%             %     \ENDIF{}
%             %     \ENDFOR
%             %     \STATE $U'\leftarrow U'\cup U_\tau $
%             % \ENDFOR
  

% 	\end{algorithmic}
% \end{algorithm}

% \begin{algorithm}[htb]
% 	\caption{Sparsification for Multiple Budget for knapsack constraint}
% 	\begin{algorithmic}[1]\label{alg:knapsack_constraint}
% 		\STATE \textbf{Input:} ground set $U$, oracle $f$, $\epsilon$, $\delta$, $B_{min}$, $B_{max}$, $w$
% 		\STATE \textbf{Output:} $U'$
%             \STATE \textbf{Initialize:} $U' \leftarrow \emptyset $ 
% 		\STATE $k \leftarrow \frac{B_{max}}{B_{min}}$ ,$ m \leftarrow \left\lceil \frac{log(k)}{log(1+\epsilon)} \right\rceil$ 
%             \FOR {$i \leftarrow 0$ to $m$} 
%                 \STATE $\tau_{i} \leftarrow B_{min}(1+ \epsilon)^i$ , $U_{\tau_{i}} \leftarrow \emptyset$
%             \ENDFOR

%             \FOR {$v \in U$} 
%                 % \STATE $\tau \leftarrow (1+ \epsilon)^i$ , $U_\tau \leftarrow \emptyset$
%                 \FOR{$i \leftarrow 0$ to $m$}
%                     \IF{$  \frac{f(U_{\tau_{i}}+v)-f(U_{\tau_{i}})}{w(v)}    \geq \frac{\delta f(U_{\tau_{i}})}{\tau_{i}}$}
%                     \STATE $U_{\tau_{i}}  \leftarrow U_{\tau_{i}} \cup v$
%                 \ENDIF{}
%                 \ENDFOR
                
%             \ENDFOR
%             \STATE $U' \leftarrow \cup_{i \in [m]} U_{\tau_{i}} $
%             % \FOR {$v \in U$} 
                
%             %     \FOR{$i \leftarrow 0$ to $m$ }
%             %         % \STATE $\tau \leftarrow (1+ \epsilon)^i$ , $U_\tau \leftarrow \emptyset$
%             %         \IF{$f(U_\tau+v)-f(U_\tau) \geq \frac{\delta f(U_\tau)}{\tau}$}
%             %         \STATE $U_\tau\leftarrow U_\tau\cup v$
%             %     \ENDIF{}
%             %     \ENDFOR
%             %     \STATE $U'\leftarrow U'\cup U_\tau $
%             % \ENDFOR
  

% 	\end{algorithmic}
% \end{algorithm}


%%% Local Variables:
%%% mode: latex
%%% TeX-master: "main.tex"
%%% End:
